\documentclass[12pt]{article}
\usepackage[utf8]{inputenc}
\usepackage[T1]{fontenc}
\usepackage{hyperref}
\title{Taivoan Language Resources for K-12 Curriculum Development}
\author{Compiled by Claude Sonnet 5 (claude-sonnet-5) for Andr\'as Kornai}
\date{}
\begin{document}
\maketitle

\section{Introduction}
Taivoan (ISO 639-3 tvx) is a dormant Formosan (Austronesian) language of Taiwan. It has no first-language speakers. SIL recognized it as a separate language in 2019. The Taivoan are not among Taiwan's officially recognized Indigenous peoples, and revival resources are limited. The search found only a few resources, listed below. No school curriculum was found.

\section{Primary Educational Resources}
No primers, textbooks or school programmes for Taivoan turned up in the search results.

\section{Linguistic and Reference Materials}
\begin{itemize}
\item Living Dictionaries, Taivoan: an online dictionary with Mandarin and English glosses, some entries with audio. Its page says the data come from a 2003 publication of Siraya and Basai linguistic materials. \url{https://livingdictionaries.app/taivoan}
\item About page of the Taivoan Living Dictionary. \url{https://livingdictionaries.app/taivoan/about}
\item Wikipedia, Taivoan language (overview and status). \url{https://en.wikipedia.org/wiki/Taivoan_language}
\end{itemize}

\section{Religious and Community Resources}
\begin{itemize}
\item Wikitongues, ``Reawakening the Taivoan Language'': a blog post on the revival effort. \url{https://wikitongues.org/blog/reawakening-the-taivoan-language/}
\item Wikitongues profile of Benson Ko-Chou Fang, who leads a project to compile a dictionary covering the four Taivoan subgroups (Tevorangh, Vogavon, Kapoa, Sia-urie). The search results do not say whether the dictionary is published. \url{https://wikitongues.org/fellows/benson-ko-chou-fang/}
\item Wikipedia, Taivoan people. \url{https://en.wikipedia.org/wiki/Taivoan}
\end{itemize}

\section{Digital and Audio Resources}
The Living Dictionaries site (above) is the only audio-bearing resource found. A Taivoan Wikimedia Incubator wiki was reported in 2019, but I did not see its URL in the search results.

\section{Curriculum Development Context}
Taivoan is dormant and thinly documented. A curriculum would have to rest on the Living Dictionaries data and the community dictionary project, and the source material should be verified with the community. The materials found are not enough for a full K-12 programme.
\end{document}